%% CO1a - Graphs - v1 - Fall 2026 - Physics 201
%%
%% Rebuild of the Fall 2025 CO1a graphs assessment on the jjc-assessment
%% template. Every figure is drawn by matplotlib inside the pycode block
%% below rather than pasted in as an image, so each figure reads its numbers
%% from the PARAMS dictionary and can be randomized per student.
%%
%% RANDOMIZE = False reproduces the Fall 2025 figures exactly.
%% RANDOMIZE = True  draws a fresh variant; every solution is recomputed
%%                   from the drawn values, so the answer key always matches.
%%
%% Build:
%%     pdflatex -interaction=nonstopmode "CO1a - Graphs - v1 - Fall 2026 - Physics 201.tex"
%%     pythontex "CO1a - Graphs - v1 - Fall 2026 - Physics 201.tex"
%%     pdflatex -interaction=nonstopmode "CO1a - Graphs - v1 - Fall 2026 - Physics 201.tex"

\documentclass[12pt]{exam}
\usepackage{jjc-assessment}

\renewcommand{\class}{PHYS 201}
\renewcommand{\examnum}{CO1a (v1)}
\renewcommand{\term}{Fall 2026}

%% Figures are produced by PythonTeX, so they do not exist on the first
%% pdflatex pass. This keeps that pass clean instead of erroring on a
%% missing file.
\newcommand{\pyfig}[2][\linewidth]{%
  \IfFileExists{#2}%
    {\includegraphics[width=#1]{#2}}%
    {\fbox{\parbox{0.9\linewidth}{\centering\small
      \texttt{#2} not built yet --- run \texttt{pythontex}, then \texttt{pdflatex} again.}}}%
}

\begin{document}

\assessmentheader

\begin{pycode}
import os
import random

import numpy as np
import matplotlib
matplotlib.use('Agg')
import matplotlib.pyplot as plt

from randassign import RandAssign
ra = RandAssign()

# ----------------------------------------------------------------- randomizing
# Flip to True to generate a fresh variant. Every number that appears in a
# figure is pulled through pick(), so randomizing is a one-line change and
# the solutions below recompute themselves from whatever was drawn.
RANDOMIZE = False

rng = random.Random()


def pick(default, choices):
    """Return `default` normally, or a random choice when RANDOMIZE is on."""
    return rng.choice(choices) if RANDOMIZE else default


FIGDIR = 'figures'
os.makedirs(FIGDIR, exist_ok=True)

plt.rcParams.update({
    'font.size': 11,
    'font.family': 'sans-serif',
    'axes.linewidth': 1.2,
    'savefig.bbox': 'tight',
    'savefig.pad_inches': 0.05,
})

# Every quantity that shows up in a figure lives here.
P = {
    'tmax':      pick(7, [6, 7, 8]),
    # Q1 graph A: flat, then speeding up.
    'A_x0':      pick(1, [0, 1, 2]),
    'A_tbreak':  pick(3, [2, 3, 4]),
    'A_xend':    pick(7, [6, 7, 8]),
    # Q1 graph B: rises to a peak, then falls back.
    'B_x0':      pick(1, [0, 1, 2]),
    'B_tpeak':   pick(4, [3, 4, 5]),
    'B_xpeak':   pick(4, [3, 4, 5]),
}


def savefig(fig, name):
    path = os.path.join(FIGDIR, name)
    fig.savefig(path)
    plt.close(fig)
    return path.replace('\\', '/')


def grid_axes(ax, face, grid, xlim, ylim, xlabel, ylabel):
    """Shared look for the graph-paper style panels in Q1."""
    ax.set_facecolor(face)
    ax.set_xlim(*xlim)
    ax.set_ylim(*ylim)
    ax.set_xticks(np.arange(xlim[0], xlim[1] + 1, 1))
    ax.set_yticks(np.arange(ylim[0], ylim[1] + 1, 1))
    ax.grid(True, color=grid, linewidth=1.0)
    ax.set_xlabel(xlabel, fontweight='bold')
    ax.set_ylabel(ylabel, fontweight='bold')
    ax.tick_params(colors='black')
    for s in ax.spines.values():
        s.set_color(grid)


# -------------------------------------------------------- Q1: position vs time
def position_A(t):
    """Constant position, then a constant-acceleration climb."""
    c = (P['A_xend'] - P['A_x0']) / (P['tmax'] - P['A_tbreak']) ** 2
    return np.where(t <= P['A_tbreak'], P['A_x0'],
                    P['A_x0'] + c * (t - P['A_tbreak']) ** 2)


def position_B(t):
    """Rise to a peak and fall back: constant negative acceleration."""
    k = (P['B_xpeak'] - P['B_x0']) / P['B_tpeak'] ** 2
    return P['B_xpeak'] - k * (t - P['B_tpeak']) ** 2


t = np.linspace(0, P['tmax'], 400)

fig, axes = plt.subplots(1, 2, figsize=(6.2, 2.5), gridspec_kw={'wspace': 0.45})
for ax, f in zip(axes, (position_A, position_B)):
    grid_axes(ax, '#e9edf7', '#3b5bc0', (0, P['tmax']), (0, 7), 't (s)', 'x (m)')
    ax.plot(t, f(t), color='black', linewidth=3, solid_capstyle='round')
fig_pos = savefig(fig, 'q1-position.pdf')


def blank_pair(name, ylabel, face, grid):
    # wspace keeps the right panel's y-label clear of the left panel.
    fig, axes = plt.subplots(1, 2, figsize=(6.2, 2.5), gridspec_kw={'wspace': 0.45})
    for ax in axes:
        grid_axes(ax, face, grid, (0, P['tmax']), (-3, 4), 't (s)', ylabel)
        ax.axhline(0, color=grid, linewidth=3)
    return savefig(fig, name)


fig_blank_v = blank_pair('q1-blank-v.pdf', 'v (m/s)', '#eaeee6', '#5a8a2b')
fig_blank_a = blank_pair('q1-blank-a.pdf', 'a (m/s$^2$)', '#fbeee2', '#d2691e')

# ------------------------------------------------- Q3: identify-the-graph panels
# Each entry: (symbol, curve, correct option letter).
#   A constant positive acceleration   B constant negative acceleration
#   C constant positive velocity       D constant negative velocity
#   E none of these
IDENTIFY = [
    ('v', 'parabola_up',   'E'),   # v curves, so a is not constant
    ('x', 'line_up',       'C'),   # straight x-t line: constant positive v
    ('a', 'line_from_zero', 'E'),  # a itself is changing
    ('x', 'parabola_down', 'B'),   # x-t parabola opening down
]
if RANDOMIZE:
    rng.shuffle(IDENTIFY)


def identify_curve(kind, u):
    """u runs 0..1 across the panel; returns values in roughly -1..1."""
    if kind == 'parabola_up':
        return 0.6 * (2 * u - 1) ** 2
    if kind == 'line_up':
        return 0.25 + 0.6 * u
    if kind == 'line_from_zero':
        return 0.85 * u
    if kind == 'parabola_down':
        return 2.4 * u * (1 - u)
    raise ValueError(kind)


# Generous wspace: the 0 and t labels sit outside the axes, so tight
# packing makes one panel's "t" collide with the next panel's "0".
fig, axes = plt.subplots(1, 4, figsize=(7.2, 1.95), gridspec_kw={'wspace': 0.75})
u = np.linspace(0, 1, 300)
for ax, (sym, kind, _ans) in zip(axes, IDENTIFY):
    ax.plot(u, identify_curve(kind, u), color='black', linewidth=1.6)
    ax.axhline(0, color='black', linewidth=1.2)
    ax.set_xlim(0, 1)
    ax.set_ylim(-1.15, 1.15)
    ax.set_xticks(np.linspace(0, 1, 5))
    ax.set_yticks(np.linspace(-1, 1, 5))
    ax.grid(True, color='#2b3fa0', linestyle='--', linewidth=0.8)
    ax.set_xticklabels([])
    ax.set_yticklabels([])
    ax.tick_params(length=0)
    for side in ('top', 'right'):
        ax.spines[side].set_visible(False)
    for side in ('left', 'bottom'):
        ax.spines[side].set_color('#2b3fa0')
    ax.text(-0.06, 0.96, '$%s$' % sym, transform=ax.transAxes,
            fontsize=15, ha='right', va='bottom')
    ax.text(1.02, 0.5, '$t$', transform=ax.transAxes,
            fontsize=15, ha='left', va='center')
    ax.text(-0.05, 0.5, '0', transform=ax.transAxes,
            fontsize=13, ha='right', va='center')
fig_identify = savefig(fig, 'q3-identify.pdf')

# ------------------------------------------------------------ Q5: strobe tape
# Sphere position (m) at each 1-second flash.
STROBE = {
    'A': pick([0, 2, 4, 6, 8, 10], [[0, 2, 4, 6, 8, 10], [0, 1.5, 3, 4.5, 6, 7.5]]),
    'B': pick([0, 0.3, 1, 2, 4, 8], [[0, 0.3, 1, 2, 4, 8], [0, 0.5, 1.5, 3, 5, 8]]),
    'C': pick([0, 3, 6, 9],        [[0, 3, 6, 9], [0, 2.5, 5, 7.5]]),
    'D': pick([1, 6, 9, 10],       [[1, 6, 9, 10], [1, 5, 8, 10]]),
}
STROBE_WINDOW = pick(3, [3, 4])   # rank over the first this-many seconds

fig, axes = plt.subplots(len(STROBE), 1, figsize=(6.4, 3.1))
for ax, (label, xs) in zip(axes, STROBE.items()):
    ax.annotate('', xy=(10.6, 0), xytext=(-0.3, 0),
                arrowprops=dict(arrowstyle='-|>', color='black', linewidth=1.2))
    for m in range(11):
        ax.plot([m, m], [-0.13, 0], color='black', linewidth=1.0)
        ax.text(m, -0.62, str(m), ha='center', va='center', fontsize=9)
    ax.text(11.0, -0.62, 'm', ha='center', va='center', fontsize=9)
    ax.plot(xs, [0.30] * len(xs), 'o', markersize=11,
            markerfacecolor='#c9c9c9', markeredgecolor='black', linestyle='none')
    ax.text(-1.0, 0.0, label, fontsize=15, fontweight='bold', ha='center', va='center')
    ax.set_xlim(-1.6, 11.6)
    ax.set_ylim(-1.0, 0.75)
    ax.axis('off')
fig_strobe = savefig(fig, 'q5-strobe.pdf')

strobe_avg = {}
for label, xs in STROBE.items():
    n = min(STROBE_WINDOW, len(xs) - 1)      # flashes are 1 s apart
    strobe_avg[label] = abs(xs[n] - xs[0]) / n
strobe_rank = sorted(strobe_avg, key=lambda k: -strobe_avg[k])

# Group equal speeds so the key shows ties as "C = D > A" rather than "C > D > A".
strobe_groups = []
for label in strobe_rank:
    if strobe_groups and abs(strobe_avg[label] - strobe_avg[strobe_groups[-1][0]]) < 1e-9:
        strobe_groups[-1].append(label)
    else:
        strobe_groups.append([label])
strobe_order = ' > '.join(' = '.join(g) for g in strobe_groups)

# --------------------------------------------------------- Q6: scooter v-t plot
HEADSTART = pick(2, [2, 3])
BLACK_V = pick([(0, -1), (5, 4), (10, 4)],
               [[(0, -1), (5, 4), (10, 4)], [(0, -1), (4, 4), (10, 4)]])
RED_V = pick([(HEADSTART, 0), (3, 3), (6, 3), (10, 5)],
             [[(HEADSTART, 0), (3, 3), (6, 3), (10, 5)],
              [(HEADSTART, 0), (4, 3), (7, 3), (10, 4)]])
TFINAL = pick(10, [10, 12])


def velocity(profile, tq, for_plot=False):
    """Piecewise-linear v(t).

    Before the profile starts the scooter is stationary, so v is 0 for the
    position integral. For plotting we return NaN instead, which matches the
    original figure: the red trace simply begins at its start time rather than
    running along the axis from t=0.
    """
    ts = [p[0] for p in profile]
    vs = [p[1] for p in profile]
    before = np.nan if for_plot else 0.0
    return np.where(tq < ts[0], before, np.interp(tq, ts, vs))


tt = np.linspace(0, TFINAL, 4001)
v_black = velocity(BLACK_V, tt)
v_red = velocity(RED_V, tt)

fig, ax = plt.subplots(figsize=(6.4, 3.4))
ax.plot(tt, velocity(BLACK_V, tt, for_plot=True), color='black', linewidth=2.6)
ax.plot(tt, velocity(RED_V, tt, for_plot=True), color='red', linewidth=2.6)
ax.axhline(0, color='black', linewidth=1.0)
ax.set_xlim(-0.2, TFINAL)
ax.set_ylim(-5, 5)
ax.set_xticks(range(0, TFINAL + 1))
ax.set_yticks(range(-5, 6))
ax.grid(True, color='#dddddd', linewidth=0.8)
ax.set_xlabel('Time (s)')
ax.set_ylabel('Velocity (m/s)')
for side in ('top', 'right'):
    ax.spines[side].set_visible(False)
fig_scooter = savefig(fig, 'q6-scooters.pdf')

# Positions come from integrating the velocity curves.
x_black = np.concatenate(([0.0], np.cumsum(np.diff(tt) * (v_black[:-1] + v_black[1:]) / 2)))
x_red = np.concatenate(([0.0], np.cumsum(np.diff(tt) * (v_red[:-1] + v_red[1:]) / 2)))

lead = 'black' if x_black[-1] > x_red[-1] else 'red'

# Where do the two scooters meet again? Look for sign changes in the gap.
gap = x_black - x_red
meet = []
for i in range(1, len(tt)):
    if gap[i - 1] == 0 and tt[i - 1] > 1e-9:
        meet.append(tt[i - 1])
    elif gap[i - 1] * gap[i] < 0:
        frac = gap[i - 1] / (gap[i - 1] - gap[i])
        meet.append(tt[i - 1] + frac * (tt[i] - tt[i - 1]))
# Collapse near-duplicates from the fine sampling.
meet_clean = []
for m in meet:
    if m > 1e-6 and not any(abs(m - k) < 0.05 for k in meet_clean):
        meet_clean.append(m)

# ------------------------------------------------------------------- solutions
cA = (P['A_xend'] - P['A_x0']) / (P['tmax'] - P['A_tbreak']) ** 2
kB = (P['B_xpeak'] - P['B_x0']) / P['B_tpeak'] ** 2

ra.addsoln(
    r'\textbf{Q1 graph A:} $v=0$ until $t=\SI{%.3g}{\second}$, then rises linearly to '
    r'$\SI{%.3g}{\meter\per\second}$ at $t=\SI{%.3g}{\second}$; '
    r'$a=0$ then constant $\SI{%.3g}{\meter\per\second\squared}$.'
    % (P['A_tbreak'], 2 * cA * (P['tmax'] - P['A_tbreak']), P['tmax'], 2 * cA))

ra.addsoln(
    r'\textbf{Q1 graph B:} $v$ falls linearly from $\SI{%.3g}{\meter\per\second}$ at $t=0$, '
    r'through zero at $t=\SI{%.3g}{\second}$, to $\SI{%.3g}{\meter\per\second}$ at $t=\SI{%.3g}{\second}$; '
    r'$a$ constant at $\SI{%.3g}{\meter\per\second\squared}$.'
    % (2 * kB * P['B_tpeak'], P['B_tpeak'],
       2 * kB * (P['B_tpeak'] - P['tmax']), P['tmax'], -2 * kB))

ra.addsoln(r'\textbf{Q3:} ' + ', '.join(
    '%d.~%s' % (i + 1, ans) for i, (_s, _k, ans) in enumerate(IDENTIFY)))

ra.addsoln(r'\textbf{Q5:} average speeds over the first \SI{%d}{\second}: '
           % STROBE_WINDOW
           + ', '.join(r'%s $=\SI{%.2f}{\meter\per\second}$' % (k, strobe_avg[k])
                       for k in strobe_rank)
           + r'. Ranking greatest to least: ' + strobe_order + '.')

ra.addsoln(r'\textbf{Q6a:} the %s scooter leads at $t=\SI{%d}{\second}$ '
           r'($\SI{%.1f}{\meter}$ vs $\SI{%.1f}{\meter}$).'
           % (lead, TFINAL, max(x_black[-1], x_red[-1]), min(x_black[-1], x_red[-1])))

ra.addsoln(r'\textbf{Q6b:} same position at $t=$ '
           + ', '.join(r'\SI{%.2f}{\second}' % m for m in meet_clean) + '.')
\end{pycode}

\begin{questions}

%% ---------------------------------------------------------------- question 1
\question For the following position vs.\ time graphs, draw the corresponding
velocity vs.\ time and acceleration vs.\ time graphs in the graphs provided below.

\begin{center}
\pyfig[5.2in]{figures/q1-position.pdf}
\end{center}

\begin{center}
\pyfig[5.2in]{figures/q1-blank-v.pdf}
\end{center}

\begin{center}
\pyfig[5.2in]{figures/q1-blank-a.pdf}
\end{center}

%% ---------------------------------------------------------------- question 2
\question Explain the overall process for drawing the velocity and acceleration
graphs. (You do not need to explain each graph individually.)

\vspace{1.75in}

\newpage

%% ---------------------------------------------------------------- question 3
\question Identify the acceleration or velocity associated with each of these
graphs from the following options:

\begin{center}
\begin{tabular}{ll}
A. Constant positive acceleration & B. Constant negative acceleration \\
C. Constant positive velocity     & D. Constant negative velocity \\
E. None of these                  &
\end{tabular}
\end{center}

\begin{center}
\pyfig{figures/q3-identify.pdf}
\end{center}

\begin{center}
\begin{tabular}{cccc}
1. \makebox[1.2in]{\hrulefill} & 2. \makebox[1.2in]{\hrulefill} &
3. \makebox[1.2in]{\hrulefill} & 4. \makebox[1.2in]{\hrulefill} \\
\end{tabular}
\end{center}

%% ---------------------------------------------------------------- question 4
\question Explain your choices.

\vspace{2.25in}

\newpage

%% ---------------------------------------------------------------- question 5
\question In each case below, a sphere moves from left to right beside a tape
marked in meters. A strobe lamp flashes once every second, and each circle
records where the sphere was at one flash. The spheres are not all shown for
the same total time.

\begin{center}
\pyfig[6.2in]{figures/q5-strobe.pdf}
\end{center}

Rank the magnitude of the average velocity of each sphere over the first
\py{STROBE_WINDOW} seconds, from greatest to least. Note any ties.

\begin{center}
\begin{tabular}{cccc}
\makebox[0.9in]{\hrulefill} & \makebox[0.9in]{\hrulefill} &
\makebox[0.9in]{\hrulefill} & \makebox[0.9in]{\hrulefill} \\
greatest & & & least \\
\end{tabular}
\end{center}

Explain your reasoning.

\vspace{1.75in}

\newpage

%% ---------------------------------------------------------------- question 6
\question Two kids are racing on their scooters. They both start at the same
location, but the kid with the black scooter (represented by the black line)
gets a \py{HEADSTART}-second head start over the kid with the red scooter
(represented by the red line). A velocity vs.\ time graph is shown below for the
two scooters.

\begin{center}
\pyfig[6.2in]{figures/q6-scooters.pdf}
\end{center}

\begin{parts}
\part After \py{TFINAL} seconds, which scooter is furthest from the starting
point?

\vspace{1.25in}

\part Identify any points (after $t=0$) where the two scooters are at the same
position.

\vspace{1.25in}
\end{parts}

\end{questions}

\end{document}
